\documentclass[12pt,letterpaper, onecolumn]{exam}
\usepackage{amsmath}
\usepackage{amssymb}
\usepackage[lmargin=71pt, tmargin=1.2in]{geometry}
\usepackage{xcolor}
\usepackage[brazil]{babel}
\usepackage{natbib}
\usepackage{enumitem}
\usepackage{booktabs}
\usepackage[hidelinks]{hyperref}
\urlstyle{same}

\renewenvironment{solution}
{\par\noindent\textbf{R:}\enspace\color{blue}}
{\par\color{black}}

% \chead{\hline} % Un-comment to draw line below header
\thispagestyle{empty}

\begin{document}

\begingroup
    \centering
    \LARGE Lista 10 - Econometria I - 2026\\[0.5em]
    %\large \today\\[0.5em]
    \large Prof: André Luiz Pereira Mancha \par
    \large Monitor: Tuffy Licciardi Issa \par
    %\large Roll Number\par
    %\large Class/Batch/Section\par
\endgroup
\rule{\textwidth}{0.4pt}
\pointsdroppedatright
\printanswers
\renewcommand{\solutiontitle}{\noindent\textbf{Ans:}\enspace}
\newcommand{\qstar}{\hspace{-0.6em}\textsuperscript{*}\hspace{0.2em}}

\begin{questions}

\question\qstar (8.5, \citep{wooldridge2016}) A variável \textit{smokes} é uma variável binária igual a um se a pessoa fuma, e zero caso contrário. Utilizando os dados contidos no arquivo \texttt{SMOKE}, estimamos um modelo de probabilidade linear de \textit{smokes}:

\[
\begin{aligned}
\widehat{smokes} =\;& 0{,}656 - 0{,}069\,\log(cigpric) + 0{,}012\,\log(income) - 0{,}029\,educ \\
& (0{,}855)\hspace{1.0cm}(0{,}204)\hspace{1.0cm}(0{,}026)\hspace{1.0cm}(0{,}006) \\
& [0{,}856]\hspace{1.0cm}[0{,}207]\hspace{1.0cm}[0{,}026]\hspace{1.0cm}[0{,}006] \\
& +\; 0{,}020\,age - 0{,}00026\,age^2 - 0{,}101\,restaurn - 0{,}026\,white \\
& \hspace{0.55cm}(0{,}006)\hspace{1.0cm}(0{,}00006)\hspace{1.0cm}(0{,}039)\hspace{1.0cm}(0{,}052) \\
& \hspace{0.55cm}[0{,}005]\hspace{1.0cm}[0{,}00006]\hspace{1.0cm}[0{,}038]\hspace{1.0cm}[0{,}050]
\end{aligned}
\]
\[
n = 807,\quad R^2 = 0{,}062.
\]

A variável \textit{white} é igual a um se a pessoa for branca, e zero caso contrário; as outras variáveis independentes foram definidas no Exemplo 8.7. Tanto os erros padrão usuais quanto os robustos em relação à heteroscedasticidade estão informados.

\begin{enumerate}[label=(\roman*)]
    \item Existe alguma diferença importante entre os dois conjuntos de erros padrão?
    \item Mantendo os outros fatores fixos, se a educação aumentar em quatro anos, o que acontece com a probabilidade estimada de fumar?
    \item Em que ponto mais um ano de idade reduz a probabilidade de fumar?
    \item Interprete o coeficiente da variável binária \textit{restaurn} (uma variável \textit{dummy} igual a um se a pessoa viver em um Estado em que há restrições ao fumo em restaurantes).
    \item A pessoa número 206 no conjunto de dados tem as seguintes características: \(cigpric = 67{,}44\), \(income = 6.500\), \(educ = 16\), \(age = 77\), \(restaurn = 0\), \(white = 0\) e \(smokes = 0\). Calcule a probabilidade de fumar dessa pessoa e comente o resultado.
\end{enumerate}

\vspace{0.7cm}

\question (8.7, \citep{wooldridge2016}) Considere um modelo em nível de empregado,
\[
y_{i,e} = \beta_0 + \beta_1 x_{i,e,1} + \beta_2 x_{i,e,2} + \cdots + \beta_k x_{i,e,k} + f_i + v_{i,e},
\]
em que a variável não observada \(f_i\) é um ``efeito da empresa'' para cada empregado da empresa \(i\). O termo de erro \(v_{i,e}\) é específico do empregado e da empresa. O erro de combinação é \(u_{i,e} = f_i + v_{i,e}\), como na equação (8.28).
\begin{enumerate}[label=(\roman*)]
    \item Suponha que \(\operatorname{Var}(f_i)=\sigma_f^2\), \(\operatorname{Var}(v_{i,e})=\sigma_v^2\), e que \(f_i\) e \(v_{i,e}\) sejam não correlacionadas. Demonstre que \(\operatorname{Var}(u_{i,e})=\sigma_f^2+\sigma_v^2\). Chame esse valor de \(\sigma^2\).
    \item Agora suponha que, para \(e \neq g\), \(v_{i,e}\) e \(v_{i,g}\) sejam não correlacionadas. Demonstre que \(\operatorname{Cov}(u_{i,e},u_{i,g})=\sigma_f^2\).
    \item Seja
    \[
    \bar{u}_i = m_i^{-1}\sum_{e=1}^{m_i} u_{i,e}
    \]
    a média dos erros de combinação em determinada empresa. Demonstre que
    \[
    \operatorname{Var}(\bar{u}_i)=\sigma_f^2+\frac{\sigma_v^2}{m_i}.
    \]
    \item Detalhe a relevância do item (iii) para a estimação por MQP usando dados nivelados pela média da empresa, em que o peso usado para a observação \(i\) é o tamanho da empresa.
\end{enumerate}

\vspace{0.7cm}

\question\qstar (8.1, \citep{wooldridge2016}) Quais das seguintes alternativas são consequências da heteroscedasticidade?
\begin{enumerate}[label=(\roman*)]
    \item Os estimadores MQO, \(\hat{\beta}_j\), são inconsistentes.
    \item A estatística \(F\) usual não mais tem uma distribuição \(F\).
    \item Os estimadores MQO não são mais BLUE.
\end{enumerate}

\vspace{0.7cm}

\question (8.3, \citep{wooldridge2016}) Verdadeiro ou falso? O método MQP é preferido ao MQO quando uma variável importante for omitida.

\vspace{0.7cm}

\question\qstar (8.C13, \citep{wooldridge2016}) Use os dados do arquivo \texttt{FERTIL2} para esta questão.
\begin{enumerate}[label=(\roman*)]
    \item Estime o modelo
    \[
    children = \beta_0 + \beta_1 age + \beta_2 age^2 + \beta_3 educ + \beta_4 electric + \beta_5 urban + u
    \]
    e registre os erros padrão usuais e robustos em relação à heteroscedasticidade. Os erros padrão robustos são sempre maiores do que os não robustos?
    \item Adicione três variáveis \textit{dummy} de religião e teste se elas são conjuntamente significativas. Quais são os \(p\)-valores dos testes não robustos e robustos?
    \item A partir da regressão do item (ii), obtenha os valores ajustados \(\hat{y}_i\) e os resíduos \(\hat{u}_i\). Regrida \(\hat{u}_i^2\) sobre \(\hat{y}_i\) e \(\hat{y}_i^2\), e teste a significância conjunta dos dois regressores. Conclua que a heteroscedasticidade está presente na equação para \textit{children}.
    \item Você diria que a heteroscedasticidade encontrada no item (iii) é importante na prática?
\end{enumerate}

\end{questions}

\bibliographystyle{apalike}
\bibliography{refs}
\end{document}
